\documentclass[10pt,a4paper]{amsart}
\usepackage{amsmath,amssymb,amsthm,graphicx}
% For commutative diagrams we use Xy-pic
\usepackage[all]{xy}

% We use 2cell for 2-commutative diagrams.
\xyoption{2cell}
% Theorem environments.
%
\theoremstyle{plain}
\newtheorem{theorem}[subsection]{Theorem}
\newtheorem{proposition}[subsection]{Proposition}
\newtheorem{lemma}[subsection]{Lemma}

\theoremstyle{definition}
\newtheorem{definition}[subsection]{Definition}
\newtheorem{example}[subsection]{Example}
\newtheorem{exercise}[subsection]{Exercise}
\newtheorem{situation}[subsection]{Situation}

\theoremstyle{remark}
\newtheorem{remark}[subsection]{Remark}
\newtheorem{remarks}[subsection]{Remarks}

\numberwithin{equation}{subsection}

% Macros
%
\def\lim{\mathop{\mathrm{lim}}\nolimits}
\def\colim{\mathop{\mathrm{colim}}\nolimits}
\def\Spec{\mathop{\mathrm{Spec}}}
\def\Hom{\mathop{\mathrm{Hom}}\nolimits}
\def\Ext{\mathop{\mathrm{Ext}}\nolimits}
\def\SheafHom{\mathop{\mathcal{H}\!\mathit{om}}\nolimits}
\def\SheafExt{\mathop{\mathcal{E}\!\mathit{xt}}\nolimits}
\def\Sch{\mathit{Sch}}
\def\Mor{\mathop{\mathrm{Mor}}\nolimits}
\def\Ob{\mathop{\mathrm{Ob}}\nolimits}
\def\Sh{\mathop{\mathit{Sh}}\nolimits}
\def\NL{\mathop{N\!L}\nolimits}
\def\CH{\mathop{\mathrm{CH}}\nolimits}
\def\proetale{{pro\text{-}\acute{e}tale}}
\def\etale{{\acute{e}tale}}
\def\QCoh{\mathit{QCoh}}
\def\Ker{\mathop{\mathrm{Ker}}}
\def\Im{\mathop{\mathrm{Im}}}
\def\Coker{\mathop{\mathrm{Coker}}}
\def\Coim{\mathop{\mathrm{Coim}}}

% Boxtimes
%
\DeclareMathSymbol{\boxtimes}{\mathbin}{AMSa}{"02}

%
% Macros for moduli stacks/spaces
%
\def\QCohstack{\mathcal{QC}\!\mathit{oh}}
\def\Cohstack{\mathcal{C}\!\mathit{oh}}
\def\Spacesstack{\mathcal{S}\!\mathit{paces}}
\def\Quotfunctor{\mathrm{Quot}}
\def\Hilbfunctor{\mathrm{Hilb}}
\def\Curvesstack{\mathcal{C}\!\mathit{urves}}
\def\Polarizedstack{\mathcal{P}\!\mathit{olarized}}
\def\Complexesstack{\mathcal{C}\!\mathit{omplexes}}
% \Pic is the operator that assigns to X its picard group, usage \Pic(X)
% \Picardstack_{X/B} denotes the Picard stack of X over B
% \Picardfunctor_{X/B} denotes the Picard functor of X over B
\def\Pic{\mathop{\mathrm{Pic}}\nolimits}
\def\Picardstack{\mathcal{P}\!\mathit{ic}}
\def\Picardfunctor{\mathrm{Pic}}
\def\Deformationcategory{\mathcal{D}\!\mathit{ef}}
\usepackage{cleveref}
\crefname{originallemma}{original lemma}{original lemmas}
\Crefname{originallemma}{Original lemma}{Original lemmas}
\makeatletter
\newlabel{lemma-proper-pushforward-stalk}{{0A3T}{source}{}{sourcepointer-1}{}}
\newlabel{lemma-proper-pushforward-stalk@cref}{{[originallemma][0][]0A3T}{[1][0][]source}{}{}{}}
\makeatother
\title{Base change of pushforward along a proper morphism without the finiteness hypotheses}
\author{The Stacks Project Authors}
\date{stacks-project:etale-cohomology.tex:local-snapshot-20261011; GNU FDL 1.2 or later}
\begin{document}
\maketitle
\noindent\textit{Editorial note.} All delivered mathematical text is the author's own, copied from the pinned Stacks Project chapter file etale-cohomology.tex; no statement or proof step was rewritten, reordered or removed.\par\medskip
\noindent\textit{Editorial note.} The proof is the author's complete argument up to the closing sentence, in which the verification that the stalk isomorphism is compatible with the constructed comparison map is marked as omitted. That verification is a routine tail of the argument, and the sentence is retained verbatim inside the counted proof span.\par\medskip
\noindent\textit{Editorial note.} The label lemma-proper-base-change-f-star occurs in no other staged Stacks chapter, so no cross-chapter restatement had to be weighed for this card.\par\medskip
\noindent\textit{Editorial note.} The single labelled reference occurring inside the counted spans is carried in assembly.external\_references as a non-load-bearing navigation pointer to the same chapter.\par\medskip
\noindent\textit{Editorial note.} The preamble is the chapter's own native preamble excerpt already used for this source: the Xy-pic setup (native preamble lines 14-18) and the native theorem declarations with the mathematical macros (native preamble lines 154-217).\par\medskip
\noindent\textit{Allowed standard prerequisites (editorial).}\begin{itemize}
\item Pushforward and inverse image of sheaves on small etale sites. Same chapter, the functors f\_* and g\textasciicircum{}\{-1\} on small etale sites and their adjunction. Supplies the canonical map g\textasciicircum{}\{-1\}f\_*F -> f'\_*(g')\textasciicircum{}\{-1\}F constructed at the start of the proof.
\item Geometric fibres X\_y and X'\_\{y'\} of a morphism and their identification under base change. Same chapter, the description of the geometric fibre X\_y = X x\_Y Spec(k) of a morphism at a geometric point y : Spec(k) -> Y. Supplies the identification X\_y = X'\_\{y'\} used at the end of the proof.
\item Stalks of a pushforward along a proper morphism. Same chapter, the identification of the stalk (f\_*F)\_\{y\} with the sections of F over the geometric fibre X\_y for proper f. Supplies the description of both sides on stalks, which is what the proof compares.
\end{itemize}\medskip
\ifdefined\hypertarget\hypertarget{sourcepointer-1}{}\fi
\noindent\textit{Source reference.} Original lemma 0A3T: sources/stacks-etale-cohomology/pinned-original-source.tex, source lines 18254--18254. Same-chapter navigation pointer invoked at proof lines 18299-18300, where the author identifies the stalks of the pushforward with sections over the geometric fibre. The pointer only resolves that reference for the reader; the cited lemma is not reproduced inside the counted span, and no bibliography citation is involved.\par\medskip
\begin{lemma}
\label{lemma-proper-base-change-f-star}
Let $f : X \to Y$ be a proper morphism of schemes. Let $g : Y' \to Y$
be a morphism of schemes. Set $X' = Y' \times_Y X$ with projections
$f' : X' \to Y'$ and $g' : X' \to X$. Let $\mathcal{F}$ be any sheaf on
$X_\etale$. Then $g^{-1}f_*\mathcal{F} = f'_*(g')^{-1}\mathcal{F}$.
\end{lemma}

\begin{proof}
There is a canonical map $g^{-1}f_*\mathcal{F} \to f'_*(g')^{-1}\mathcal{F}$.
Namely, it is adjoint to the map
$$
f_*\mathcal{F} \longrightarrow
g_*f'_*(g')^{-1}\mathcal{F} = f_*g'_*(g')^{-1}\mathcal{F}
$$
which is $f_*$ applied to the canonical map
$\mathcal{F} \to g'_*(g')^{-1}\mathcal{F}$. To check this map is an
isomorphism we can compute what happens on stalks.
Let $y' : \Spec(k) \to Y'$ be a geometric point with image $y$ in $Y$.
By Lemma \ref{lemma-proper-pushforward-stalk} the stalks are
$\Gamma(X'_{y'}, \mathcal{F}_{y'})$ and $\Gamma(X_y, \mathcal{F}_y)$
respectively. Here the sheaves $\mathcal{F}_y$ and $\mathcal{F}_{y'}$
are the pullbacks of $\mathcal{F}$ by the projections $X_y \to X$
and $X'_{y'} \to X$. Since $X_y = X'_{y'}$ as schemes
we see that the stalks agree. We omit the
verification that this isomorphism is compatible with our map.
\end{proof}

\end{document}
